\section{Introdução}
\label{sec-intro}

A resolução de Problemas de Otimização Combinatória apresenta desafios significativos devido à sua complexidade, que inviabiliza a obtenção de soluções ótimas por meio de algoritmos exatos em grandes instâncias, dado o crescimento exponencial do espaço de busca. Nesse contexto, as meta-heurísticas surgem como alternativas eficazes, capazes de produzir soluções de alta qualidade em tempo computacional viável. Embora não garantam soluções ótimas, sua eficiência e adaptabilidade as tornam amplamente estudadas e aprimoradas. Uma tendência atual é o desenvolvimento de abordagens híbridas, que combinam o potencial exploratório das meta-heurísticas com a precisão de métodos exatos, buscando unir o melhor de ambos os mundos para resolver problemas de grande escala de forma mais robusta e eficiente, como abordado nos estudos em \cite{Puchinger:2005}.

A Teoria dos Grafos desempenha papel central nesse campo, fornecendo modelos baseados em vértices e arestas que representam de maneira natural as relações entre entidades. Essa representação permite formular e resolver problemas clássicos de otimização, como a árvore geradora mínima, que consiste em encontrar uma subestrutura sem ciclos que conecte todos os vértices com o menor custo total. Tais problemas têm aplicações diretas em redes de transporte e telecomunicações, onde a minimização de custos é fundamental. Propriedades estruturais dos grafos, como medidas de centralidade, também auxiliam na compreensão da relevância dos vértices em uma rede. A centralidade de grau, por exemplo, indica a conectividade direta de um vértice, sendo útil na análise de redes sociais ou urbanas. Outras métricas, como proximidade e intermediação, fornecem perspectivas adicionais sobre influência e posição relativa na rede.

Entre os problemas abordados por meio da Teoria dos Grafos destaca-se para este estudo o Problema da Árvore Geradora com Número Mínimo de Vértices Branch, ou \textit{Minimum Branch Vertices Problem} (MBV). Esse problema consiste em identificar, em um grafo conexo e não direcionado, uma árvore geradora que minimize o número de vértices com grau superior a $d$, onde $d \ge 2$ . Sua relevância prática está associada à alocação otimizada de \textit{switches} WDM (Multiplexação por Divisão de Comprimento de Onda) em redes ópticas voltadas à transmissão \textit{multicast}, onde a minimização desses dispositivos, de alto custo, representa um objetivo estratégico.

O presente trabalho se baseia em estudos anteriores de \cite{silva2024ils}, \cite{lima2024ils} e \cite{oliveira:2023}, propondo uma nova abordagem construtiva para o problema de minimização do número de vértices \textit{branch} em árvores geradoras. A abordagem integra a análise da densidade do grafo e da centralidade \textit{PageRank} com um mecanismo de aleatoriedade para guiar a geração de soluções iniciais. Essa estratégia será aplicada em meta-heurísticas consolidadas na literatura, como \textit{GRASP} e \textit{Multi-Start}, a fim de avaliar seu impacto na qualidade e na eficiência das soluções obtidas para o d-MBV.

Os objetivos principais deste estudo são aprimorar a busca por soluções eficientes para o problema d-MBV e divulgar a experiência e os resultados alcançados durante a pesquisa. O trabalho está estruturado da seguinte forma: inicia-se com o embasamento teórico com os conceitos fundamentais; a metodologia descreve em detalhes o algoritmo proposto e as meta-heurísticas utilizadas; por fim, a seção de resultados e conclusões discute o desempenho das variações algorítmicas e avalia seu impacto sobre as instâncias testadas. 

